\documentclass{article}
\usepackage{amsmath}
\usepackage{amsfonts}
\usepackage{amssymb}
\usepackage{bussproofs}
\usepackage{mathtools}

\begin{document}
\section*{Task 1}
Congruence Rules:
\begin{itemize}
\item \textsc{E-OpFst}
\item \textsc{E-OpSnd} 
\item \textsc{E-IfCond}
\end{itemize}
All other rules are computation rules 

\begin{align*}
E & \coloneq & \square \\
  && E \oplus t \\ 
  && v \oplus E \\
  && \mathtt{if}(E;t_1,t_2)\\
\end{align*}

\section*{Task 2}
\begin{prooftree}
\AxiomC{}
\RightLabel{\textsc{E-OrTrueFalse}}
\UnaryInfC{$\mathtt{true} \vee \mathtt{false} \mapsto \mathtt{true}$}
\RightLabel{\textsc{E-OpFst}}
\UnaryInfC{$(\mathtt{true}\vee\mathtt{false})\wedge (\mathtt{false}\vee\mathtt{false})\mapsto \mathtt{true} \wedge (\mathtt{false} \vee \mathtt{false})$}
\RightLabel{\textsc{E-IfCond}}
\UnaryInfC{$\mathtt{if}((\mathtt{true}\vee\mathtt{false})\wedge (\mathtt{false}\vee\mathtt{false});(7*7)+15,4+4)$}
\noLine{}
\UnaryInfC{$\mapsto \mathtt{if}(\mathtt{true}\wedge (\mathtt{false}\vee\mathtt{false});(7*7)+15,4+4)$}
\end{prooftree}

\begin{itemize}
\item Every number $n\in\mathbb{Z}$, $\mathtt{true}$,$\mathtt{false}$ are values and normal forms
\item $\mathtt{if}(3;1,2)$ is in normal form, but not a value,
\item The only rules for evaluating a $\mathtt{if}$-term with the condition already a value are the rules \textsc{E-IfTrue} and \textsc{E-IfFalse}, which requires the condition to be either $\mathtt{true}$ or $\mathtt{false}$, but $3$ is neither
\end{itemize}

We claim given terms $t$ and $t^{\prime}$ with $t\mapsto t^{\prime}$, we have $size(t) \geq size(t^{\prime})$. \\
Proof by induction on the evaluation rule used to evaluate $t\mapsto t^{\prime}$

\begin{itemize}
\item \textsc{E-Op}, then $t = n_1 \oplus n_2$ for some $n_1,n_2\in\mathbb{Z}$ and binary operation $\oplus\in\{+,-,*\}$. 
  Then $t^{\prime} = n$ where $n\in\mathbb{Z}$ is the result of applying the given binary operation to the two numbers, and we have 
  \begin{align*}
    \mathtt{size}(t) = \mathtt{size}(n_1) + \mathtt{size}(n_2) + 1 = 1 + 1 + 1 \\
    \leq 1  = \mathtt{size}(n) =\mathtt{size}(t^{\prime})
  \end{align*}
\item \textsc{E-IfTrue}, then $t =\mathtt{if}(\mathtt{true};t^{\prime},t_2)$ for some term $t_2$. Then we have 
  \begin{align*}
    \mathtt{size}(t) = \mathtt{size}(\mathtt{true}) + \mathtt{size}(t^{\prime}) + \mathtt{size}(t_2) + 1 \leq \mathtt{size}(t^{\prime})
  \end{align*}
\item \textsc{E-IfFalse}. This is completely analogous to \textsc{E-IfTrue} with the places of $t^{\prime}$ and $t_2$ exchanged, and $\mathtt{true}$ replaced with $\mathtt{false}$
\item \textsc{E-OrTrueTrue}, then $t = \mathtt{true} \vee \mathtt{true}$ and $t^{\prime}$, and we have 
  \begin{align*}
    \mathtt{size}(t) = \mathtt{size}(\mathtt{true)} + \mathtt{size}(\mathtt{true}) + 1 \leq \mathtt{size}(\mathtt{true}) = \mathtt{size}(t^{\prime})
  \end{align*}
\item \textsc{E-OrTrueFalse}, \textsc{E-OrFalseTrue}, \textsc{E-OrFalseFalse}, \textsc{E-AndTrueTrue}, \textsc{E-AndTrueFalse}, \textsc{E-AndFalseTrue}, \textsc{E-AndFalseFalse} are all analogous to the \textsc{E-OrTrueTrue} case, only with $\mathtt{true}$ and $\vee$ replaced by $\mathtt{false}$ and $\wedge$ depending on the case
\item \textsc{E-OpFst}, then $t=t_1 \oplus t_2$ and $t^{\prime} = t_1^{\prime}\oplus t_2$ for some terms $t_1,t_1^{\prime},t_2$ with $t_1\mapsto t_1^{\prime}$. By induction hypothesis, we have $\mathtt{size}(t_1)\leq \mathtt{size}(t_1^{\prime})$ and this gives us 
  \begin{align*}
    \mathtt{size}(t) = \mathtt{size}(t_1) + \mathtt{size}(t_2) + 1 \leq_{\text{IH}} \mathtt{size}(t_1^{\prime}) + \mathtt{size}(t_2) + 1 = \mathtt{size}(t^{\prime})
  \end{align*}
\item \textsc{E-OpSnd}, analogous to \textsc{E-OpFst} just with the second term instead of the first term used with the induction hypothesis 
\item \textsc{E-IfCond}, then $t = \mathtt{if}(t_1;t_2,t_3)$ and $t^{\prime} = \mathtt{ifz}(t_1^{\prime};t_2,t_3)$ for some terms $t_1,t_1^{\prime},t_2,t_3$ with $t_1\mapsto t_1^{\prime}$. By induction hypothesis, we have $\mathtt{size}(t_1) \leq \mathtt{size}(t_1^{\prime})$ which gives 
  \begin{align*}
    \mathtt{size}(t) = \mathtt{size}(t_1) + \mathtt{size}(t_2) + \mathtt{size}(t_3) + 1\leq_{\text{IH}}\\
    \mathtt{size}(t_1^{\prime} + \mathtt{size}(t_2) + \mathtt{size}(t_3) + 1 = \mathtt{size}(t^{\prime})
  \end{align*}
\end{itemize}
\end{document}
